% Chapter 13 -- Tips and Troubleshooting (Part II, Reference; AN-14
% troubleshooting form: bold symptom, then alternating Cause / Solution;
% the Mac OS 8 Installation manual Ch. 6 and the 1991 guide Ch. 16,
% apple-macintosh.md Sec 5.4, 9).
% Sources (Law 1): os/milton/kmain.c (finder_surface_open: only .EXE files
% are opened; FINDER-WIN-FULL at the fifth Finder window;
% FINDER-NEW-FOLDER-FAIL with no disk window front; FINDER_DBLCLICK_TICKS;
% the Apple menu rebuilt from the disk; flair_live_finder_key and
% flair_live_tenant_key; the move refusals and the zoom-back; the PC LOAD
% LETTER lines of panic.c and kmain.c; MSG-DOS-0014 on a failed Exit);
% os/flair/tbxgate.c (TENANT-SLOT-BUSY, TENANT-CRASH);
% os/flair/finder_ops.c (the refusal reasons); os/milton/panic.c.
% GATES: test-flair-disk-windows test-flair-app-launch test-flair-finder-cmds test-flair-file-ops test-flair-apple-menu test-flair-exit-to-dos test-flair-heap test-pubs-flair
\Gates{test-flair-disk-windows, test-flair-finder-cmds,
  test-flair-file-ops, test-flair-apple-menu, test-flair-exit-to-dos,
  test-flair-heap, test-pubs-flair}
\chapter{Tips and Troubleshooting}

When something does not work as you expect, find the description that
matches what you see.  Each lists the likely causes and what to do about
each one.

\Symptom{Double-clicking an icon does not open it}
\Cause The two clicks were too slow or the mouse moved between them.
\Solution Click twice quickly, holding the mouse still.  You can also
select a folder and choose \menuitem{Open} from the \menuname{File} menu.
\Cause The icon is a document or a program whose name ends in \file{.COM}.
\Solution Only applications, whose names end in \file{.EXE}, open from the
desktop.  Open a document from within its application.  Run a \file{.COM}
program at the InitechDOS prompt (Chapter~10).
\Cause An application from the disk is already open.
\Solution Only one can be open at a time.  Quit the open application,
then open the other.
\Cause Four Finder windows are already open.
\Solution Close one of them, then open the folder again.

\Symptom{A menu command is dimmed}
\Cause The command needs something that is not there: a selected folder
for \menuitem{Open}, a Finder window in front for \menuitem{Close Window},
\menuitem{Clean Up} and \menuitem{Arrange (by Name)}, something in the
Trash for \menuitem{Empty Trash}.
\Solution Select the item or open the window, then choose the command.
\Cause The command is not available in this version of FLAIR.
\Solution Chapter~8 lists the commands that are always dimmed.

\Symptom{New Folder does nothing}
\Cause No Finder window is in front.  A new folder is made in the folder of
the active Finder window.
\Solution Open the disk or the folder in which you want the new folder,
click its window, and choose \menuitem{New Folder} again.
\Cause The folder already holds folders named \file{NEWFOLD} and
\file{NEWFOL02} to \file{NEWFOL99}, or the disk is full.
\Solution Move or throw away some of them, or make room on the disk.

\Symptom{A keyboard equivalent does nothing}
\Cause The command belongs to a menu of another application.  The keys go
only to the menus of the active application.
\Solution Click a window of the application first.
\Cause The command is dimmed.
\Solution See ``A menu command is dimmed'', above.
\Cause HELLO or NOTES is the active application.
\Solution Their menus have no keyboard equivalents that work.  Click the
desktop to use the Finder's.

\Symptom{An icon slides back when you drop it}
\Cause The folder you dropped it in already holds something of the same
name.
\Solution Throw away or move the other item first.
\Cause You dropped a folder into itself or one of its own folders, or back
into the folder it came from.
\Solution Drop it somewhere else.
\Cause The item is locked.
\Solution If you were dragging it to the Trash, an alert says so.  Unlock
it at the InitechDOS prompt with \file{ATTRIB}, described in the
\doctitle{InitechDOS Reference}.

\Symptom{The Calculator is not in the Apple menu}
\Cause The program \file{CALCULAT.EXE} is not in the \file{APPS} folder.
\Solution Copy it back from your backup copy of the data disk.  The menu
lists an accessory only while its program is on the disk.

\Symptom{An application's window disappears}
\Cause The application stopped working, and FLAIR closed it.
\Solution Open the application again.  Work you had not saved is lost:
save often.

\Symptom{Restart, Shut Down or Exit to InitechDOS does nothing}
\Cause The open application asked about unsaved work and you clicked
\texttt{Cancel}.
\Solution Choose the command again, and this time click \texttt{Save} or
\texttt{Don't Save}.

\Symptom{The screen shows ``Bad or missing Command Interpreter''}
\Cause After \menuitem{Exit to InitechDOS}, InitechDOS could not load its
command processor from the startup disk.
\Solution Restart the computer.  If it happens again, the startup disk is
damaged: make a new copy from your original disks (see \doctitle{Getting
Started with FLAIR}).

\Symptom{The screen shows ``PC LOAD LETTER''}
\Cause InitechOS stopped because of an error it could not recover from.
The words in parentheses after the message name the error:
\texttt{(FLAIR heap exhausted)} means the computer ran out of memory for
FLAIR; \texttt{(FLAIR needs a 640x480 framebuffer)} means the display cannot
show FLAIR; \texttt{(exception} \textit{nn}\texttt{)} means a program
error.
\Solution Write down the whole message, then restart the computer.
\textit{Getting Started with FLAIR} lists what your computer needs to run
FLAIR.  If the message appears again, call the Initech Help Desk
(extension 2504) with the message you wrote down.
